\ifdefined\gkver\else\def\gkver{student}\fi
\documentclass[\gkver]{gaokaozhenti}
\begin{document}

\chapter{2009年上海}

\begin{enumerate}
\item
%% number: 1
%% paperName: 2009年上海高考第 1 题 4 分
%% typeId: 1
%% type: 单选题
%% chapter: 原子和原子核
%% point: 天然放射性
%% method: 无
%% score: 4
%% degree: 650
%% duplicateId: 0
%% body:
放射性元素衰变时放出三种射线，按穿透能力由强到弱的排列顺序是\xzanswer{B}
\fourchoices[answer=B]
{$\alpha$ 射线，$\beta$ 射线，$\gamma$ 射线}
{$\gamma$ 射线，$\beta$ 射线，$\alpha$ 射线}
{$\gamma$ 射线，$\alpha$ 射线，$\beta$ 射线}
{$\beta$ 射线，$\alpha$ 射线，$\gamma$ 射线}
%% answer: B
%% memo:
\memoanswer{
三种射线的穿透能力 $\gamma$ 射线最强、$\alpha$ 射线最弱，故 B 正确。
}

\item
%% number: 2
%% paperName: 2009年上海高考第 2 题 4 分
%% typeId: 1
%% type: 单选题
%% chapter: 气体性质
%% point: 物体的内能
%% method: 无
%% score: 4
%% degree: 650
%% duplicateId: 0
%% body:
气体内能是所有气体分子热运动动能和势能的总和，其大小与气体的状态有关，分子热运动的平均动能与分子间势能分别取决于气体的\xzanswer{A}
\fourchoices[answer=A]
{温度和体积}
{体积和压强}
{温度和压强}
{压强和温度}
%% answer: A
%% memo:
\memoanswer{
分子热运动的平均动能由温度决定，分子间的势能由气体的体积决定，故 A 正确。
}

\item
%% number: 3
%% paperName: 2009年上海高考第 3 题 4 分
%% typeId: 1
%% type: 单选题
%% chapter: 电场
%% point: 电场中的图象
%% method: 无
%% score: 4
%% degree: 450
%% duplicateId: 0
%% body:
两带电量分别为 $q$ 和 $- q$ 的点电荷放在 $x$ 轴上，相距为 $L$，能正确反映两电荷连线上场强大小 $E$ 与 $x$ 关系的是图\xzanswer{A}
\fourchoices[answer=A, ispicture=true, h=2.0cm]
{figs/03a.png}
{figs/03b.png}
{figs/03c.png}
{figs/03d.png}
%% answer: A
%% memo:
\memoanswer{
本题可以用排除法得出正确结果。两个点电荷的电荷量分别为 $q$ 和 $-q$，在二者连线上产生的场强应该是对称的，故 C、D 错误；在两个点电荷连线的中点，$q$ 和 $-q$ 产生的场强的方向相同，合场强不可能为零，故 B 错误，只有 A 正确。
}

\item
%% number: 4
%% paperName: 2009年上海高考第 4 题 4 分
%% typeId: 1
%% type: 单选题
%% chapter: 机械振动
%% point: 单摆
%% method: 无
%% score: 4
%% degree: 650
%% duplicateId: 0
%% body:
做简谐振动的单摆摆长不变，若摆球质量增加为原来的 4 倍，摆球经过平衡位置时速度减小为原来的 $ \frac{1}{2} $，则单摆振动的\xzanswer{C}
\fourchoices[answer=C]
{频率、振幅都不变}
{频率、振幅都改变}
{频率不变、振幅改变}
{频率改变、振幅不变}
%% answer: C
%% memo:
\memoanswer{
单摆的周期公式为 $T = 2\pi\sqrt{\frac{l}{g}}$，所以单摆的周期与摆球质量和摆球在平衡位置的速度无关，摆球的周期不变，频率等于周期的倒数，也不变。摆球在平衡位置的动能为 $E_{k} = \frac{1}{2}mv^{2}$，若摆球质量变为原来的 4 倍，速度变为原来的 $\frac{1}{2}$，则摆球的动能不变，总能量不变；若单摆振动的振幅不变，由于摆球质量变为原来的 4 倍，摆球的势能将会增大，所以摆球的振幅要减小，故 C 正确。
}

\item
%% number: 5
%% paperName: 2009年上海高考第 5 题 4 分
%% typeId: 1
%% type: 单选题
%% chapter: 功和能
%% point: 动能定理
%% method: 无
%% score: 4
%% degree: 450
%% duplicateId: 0
%% body:
小球由地面竖直上抛，上升的最大高度为 $H$，设所受阻力大小恒定，地面为零势能面。在上升至离地高度 $h$ 处，小球的动能是势能的两倍，在下落至离地高度 $h$ 处，小球的势能是动能的两倍，则 $h$ 等于\xzanswer{D}
\fourchoices[answer=D]
{$\frac{H}{9}$}
{$\frac{2H}{9}$}
{$\frac{3H}{9}$}
{$\frac{4H}{9}$}
%% answer: D
%% memo:
\memoanswer{
小球上升至 $h$ 处时，动能为 $2mgh$，设小球上抛的初速度为 $v$，则有 $-(f + mg)h = 2mgh - \frac{1}{2}mv^{2}$；由小球上升的最大高度为 $H$ 得 $-(f + mg)H = 0 - \frac{1}{2}mv^{2}$；小球下落至离地 $h$ 处时，动能为 $\frac{1}{2}mgh$，有 $(mg - f)(H - h) = \frac{1}{2}mgh - 0$。联立以上各式解得 $h = \frac{4}{9}H$，故 D 正确。
}

\item
%% number: 6
%% paperName: 2009年上海高考第 6 题 5 分
%% typeId: 2
%% type: 多选题
%% chapter: 光学
%% point: 光电效应
%% method: 无
%% score: 5
%% degree: 650
%% duplicateId: 0
%% body:
光电效应的实验结论是：对于某种金属\xzanswer{AD}
\fourchoices[answer=AD]
{无论光强多强，只要光的频率小于极限频率就不能产生光电效应}
{无论光的频率多低，只要光照时间足够长就能产生光电效应}
{超过极限频率的入射光强度越弱，所产生的光电子的最大初动能就越小}
{超过极限频率的入射光频率越高，所产生的光电子的最大初动能就越大}
%% answer: AD
%% memo:
\memoanswer{
能否发生光电效应与入射光的频率和极限频率的大小有关，与光照时间无关，故 A 正确、B 错误；由公式 $E_{k} = h\nu - W_{0}$ 可以看出，入射光的频率越大，所产生光电子的最大初动能就越大，与入射光强度无关，故 C 错误、D 正确。
}

\item
%% number: 7
%% paperName: 2009年上海高考第 7 题 5 分
%% typeId: 2
%% type: 多选题
%% chapter: 电场
%% point: 等势面
%% method: 无
%% score: 5
%% degree: 450
%% duplicateId: 0
%% body:
位于 A、B 处的两个带有不等量负电的点电荷在平面内电势分布如图所示，图中实线表示等势线，则\xzanswer{CD}
\onepicture[h=3.2cm]{figs/07.png}
\fourchoices[answer=CD]
{a 点和 b 点的电场强度相同}
{正电荷从 c 点移到 d 点，电场力做正功}
{负电荷从 a 点移到 c 点，电场力做正功}
{正电荷从 e 点沿图中虚线移到 f 点，电势能先减小后增大}
%% answer: CD
%% memo:
\memoanswer{
a、b 两点的电场强度方向不同，所以场强一定不同，故 A 错误；c 点电势比 d 点电势低，将正电荷由低电势点移到高电势点，电场力做负功，故 B 错误；a 点电势比 c 点电势低，将负电荷由低电势点移到高电势点，电场力做正功，故 C 正确；从 e 点到 f 点，电势先降低后升高，正电荷从 e 点到 f 点电场力先做正功、后做负功，电势能先减小后增大，故 D 正确。
}

\item
%% number: 8
%% paperName: 2009年上海高考第 8 题 5 分
%% typeId: 2
%% type: 多选题
%% chapter: 曲线运动
%% point: 万有引力定律
%% method: 无
%% score: 5
%% degree: 650
%% duplicateId: 0
%% body:
牛顿以天体之间普遍存在着引力为依据，运用严密的逻辑推理，建立了万有引力定律。在创建万有引力定律的过程中，牛顿\xzanswer{ABC}
\fourchoices[answer=ABC]
{接受了胡克等科学家关于“吸引力与两中心距离的平方成反比”的猜想}
{根据地球上一切物体都以相同加速度下落的事实，得出物体受地球的引力与其质量成正比，即 $F \propto m$ 的结论}
{根据 $F \propto m$ 和牛顿第三定律，分析了地月间的引力关系，进而得出 $F \propto m_{1}m_{2}$}
{根据大量实验数据得出了比例系数 $G$ 的大小}
%% answer: ABC
%% memo:
\memoanswer{
万有引力定律 $F = G\frac{m_{1}m_{2}}{r^{2}}$ 是牛顿猜想出来的，通过这个式子并结合物理学史，可以得出 A、B、C 正确。引力常量是卡文迪许测出的，故 D 错误。
}

\item
%% number: 9
%% paperName: 2009年上海高考第 9 题 5 分
%% typeId: 2
%% type: 多选题
%% chapter: 气体性质
%% point: 液柱移动定性分析
%% method: 无
%% score: 5
%% degree: 450
%% duplicateId: 0
%% body:
如图为竖直放置的上细下粗的密闭细管，水银柱将气体分隔成 A、B 两部分，初始温度相同。使 A、B 升高相同温度达到稳定后，体积变化量为 $\Delta V_{A}$、$\Delta V_{B}$，压强变化量为 $\Delta p_{A}$、$\Delta p_{B}$，对液面压力的变化量为 $\Delta F_{A}$、$\Delta F_{B}$，则\xzanswer{AC}
\onepicture[h=3.4cm]{\includesvg{figs/09}}
\fourchoices[answer=AC]
{水银柱向上移动了一段距离}
{$\Delta V_{A} < \Delta V_{B}$}
{$\Delta p_{A} > \Delta p_{B}$}
{$\Delta F_{A} = \Delta F_{B}$}
%% answer: AC
%% memo:
\memoanswer{
以粗管中的气体为研究对象，由 $\frac{pV}{T}=$ 常量可得，若 $V$ 不变，$T$ 变大时，$p$ 也要变大；由于水银柱下面气体压强比上面气体压强大，所以在上、下方气体升高相同温度后，水银柱下面受到的气体压强增加得更大，就打破了原有的平衡，水银柱向上运动，故 A 正确。水银柱的体积不会改变，所以气体的总体积不变，故 B 错误。水银柱上升，粗管中的一部分水银到了细管中，水银柱的总长度要变长，上下气体的压强差变大，即 $\Delta p_{A}>\Delta p_{B}$，故 C 正确。由于上下气体的压强差发生了变化，再由压力 $F=pS$ 易得，$\Delta F_{A}$ 和 $\Delta F_{B}$ 发生变化，但不相等，故 D 错误。
}

\item
%% number: 10
%% paperName: 2009年上海高考第 10 题 4 分
%% typeId: 3
%% type: 填空题
%% chapter: 光学
%% point: 光的干涉
%% method: 无
%% score: 4
%% degree: 650
%% duplicateId: 0
%% body:
如图为双缝干涉的实验示意图，若要使干涉条纹的间距变大可改用波长更\tkanswer[1.0cm]{长}（填长、短）的单色光，或是使双缝与光屏间的距离\tkanswer[1.0cm]{增大}（填增大、减小）。
\onepicture[w=5.0cm, align=right]{\includesvg{figs/10}}
%% answer: 长，增大
%% memo:
\memoanswer{
由干涉条纹间距 $\Delta x = \frac{l}{d}\lambda$ 可得，要想增大干涉条纹间距，可以改用波长较大的单色光或者增大双缝与光屏间的距离。
}

\item
%% number: 11
%% paperName: 2009年上海高考第 11 题 4 分
%% typeId: 3
%% type: 填空题
%% chapter: 电路
%% point: 逻辑电路
%% method: 无
%% score: 4
%% degree: 650
%% duplicateId: 0
%% body:
如图为某报警装置示意图，该报警装置在一扇门、两扇窗上各装有一个联动开关，门、窗未关上时，开关不闭合，只要有一个开关未闭合，报警器就会报警。该报警装置中用了两个串联的逻辑电路，虚线框甲内应选用\tkanswer[0.8cm]{或}门电路，虚线框乙内应选用\tkanswer[0.8cm]{或}门电路（填与、非、或）。
\onepicture[w=5.6cm, align=right]{\includesvg{figs/11}}
%% answer: 或，或
%% memo:
\memoanswer{
题目中要求“只要有一个开关未闭合，报警器就会报警”，所以虚线框内都应该用“或”门电路。
}

\item
%% number: 12
%% paperName: 2009年上海高考第 12 题 4 分
%% typeId: 3
%% type: 填空题
%% chapter: 机械波
%% point: 波的图象
%% method: 无
%% score: 4
%% degree: 650
%% duplicateId: 0
%% body:
弹性绳沿 $x$ 轴放置，左端位于坐标原点，用手握住绳的左端，当 $t = 0$ 时使其开始沿 $y$ 轴做振幅为 $8\Ucm$ 的简谐振动，在 $t = 0.25\Us$ 时，绳上形成如图所示的波形，则该波的波速为\tkanswer[1.0cm]{20}$\Ucms$，$t =$\tkanswer[1.2cm]{2.75} 时，位于 $x_{2} = 45\Ucm$ 的质点 N 恰好第一次沿 $y$ 轴正向通过平衡位置。
\onepicture[w=5.6cm, align=right]{\includesvg{figs/12}}
%% answer: 20，2.75
%% memo:
\memoanswer{
由图象可得该波的波速为 $v = \frac{5}{0.25}\Ucms = 20\Ucms$。当简谐波传到 N 点时，N 点向下移动，再经过半个周期 N 点才第一次通过平衡位置向上运动。该简谐波周期为 $1\Us$，传到 N 点的时间为 $t_{1} = \frac{45}{20}\Us = 2.25\Us$，总时间 $t = 2.25\Us + 0.5\Us = 2.75\Us$。
}

\item
%% number: 13
%% paperName: 2009年上海高考第 13 题 4 分
%% typeId: 3
%% type: 填空题
%% chapter: 电磁感应
%% point: 二级电磁感应
%% method: 无
%% score: 4
%% degree: 250
%% duplicateId: 0
%% body:
如图，金属棒 ab 置于水平放置的 U 形光滑导轨上，在 ef 右侧存在有界匀强磁场 $B$，磁场方向垂直导轨平面向下，在 ef 左侧的无磁场区域 cdef 内有一半径很小的金属圆环 L，圆环与导轨在同一平面内。当金属棒 ab 在水平恒力 $F$ 作用下从磁场左边界 ef 处由静止开始向右运动后，圆环 L 有\tkanswer[1.0cm]{收缩}（填收缩、扩张）趋势，圆环内产生的感应电流\tkanswer[1.0cm]{变小}（填变大、变小、不变）。
\onepicture[w=4.2cm, align=right]{\includesvg{figs/13}}
%% answer: 收缩，变小
%% memo:
\memoanswer{
当金属棒向右移动时，导线框 abcd 中会有感应电流产生，线圈中磁通量增加，由楞次定律知，圆环有收缩的趋势。导体棒 ab 在水平方向除了受到力 $F$ 作用之外，还受到水平向左的安培力，随着 ab 速度越来越大，导体棒的加速度越来越小，导线框 abcd 中的电流变化得越来越慢，所以圆环中的电流变小。
}

\item
%% number: 14
%% paperName: 2009年上海高考第 14 题 4 分
%% typeId: 3
%% type: 填空题
%% chapter: 电路
%% point: 闭合电路欧姆定律
%% method: 无
%% score: 4
%% degree: 450
%% duplicateId: 0
%% body:
图示电路中，$R_{1} = 12\UO$，$R_{2} = 6\UO$，滑动变阻器 $R_{3}$ 上标有“$20\UO$ $2\UA$”字样，理想电压表的量程有 $0 \sim 3\UV$ 和 $0 \sim 15\UV$ 两档，理想电流表的量程有 $0 \sim 0.6\UA$ 和 $0 \sim 3\UA$ 两档。闭合电键 S，将滑片 P 从最左端向右移动到某位置时，电压表、电流表示数分别为 $2.5\UV$ 和 $0.3\UA$；继续向右移动滑片 P 到另一位置，电压表指针指在满偏的 $1/3$，电流表指针指在满偏的 $1/4$，则此时电流表示数为\tkanswer[1.0cm]{0.15}$\UA$，该电源的电动势为\tkanswer[1.0cm]{7.5}$\UV$。
\onepicture[w=5.2cm, align=right]{\includesvg{figs/14}}
%% answer: 0.15，7.5
%% memo:
\memoanswer{
滑动变阻器滑片从左向右滑动，滑动变阻器的阻值变大，电流表的读数变小，电压表的读数变大，由此可得电压表的量程为 $15\UV$，电流表的量程为 $0.6\UA$，所以电流表的读数为 $0.15\UA$，电压表的读数为 $5\UV$。由串并联知识知，第一次电路中的总电流为 $0.9\UA$，第二次电路中的总电流为 $0.45\UA$。由闭合电路的欧姆定律可得

$I_{1}\left(\frac{R_{1}R_{2}}{R_{1}+R_{2}} + R_{31} + r\right) = I_{2}\left(\frac{R_{1}R_{2}}{R_{1}+R_{2}} + R_{32} + r\right)$，

代入数据得 $0.9 \times \left(4 + \frac{2.5}{0.9} + r\right) = 0.45 \times \left(4 + \frac{5}{0.45} + r\right)$，解得 $r = \frac{14}{9}\UO$，所以电动势 $E = 0.45 \times \left(4 + \frac{5}{0.45} + \frac{14}{9}\right)\UV = 7.5\UV$。
}

\item
%% number: 15
%% paperName: 2009年上海高考第 15 题 8 分
%% typeId: 4
%% type: 实验题
%% chapter: 电路
%% point: 多用表的使用
%% method: 无
%% score: 8
%% degree: 650
%% duplicateId: 0
%% body:
\begin{enumerate}
	\item
	用多用表的欧姆档测量阻值约为几十 $\UkO$ 的电阻 $R_{x}$，以下给出的是可能的操作步骤，其中 S 为选择开关，P 为欧姆档调零旋钮，把你认为正确的步骤前的字母按合理的顺序填写在下面的横线上。

	\onepicture[w=7.0cm, align=right]{figs/15.png}

	a．将两表笔短接，调节 P 使指针对准刻度盘上欧姆档的零刻度，断开两表笔

	b．将两表笔分别连接到被测电阻的两端，读出 $R_{x}$ 的阻值后，断开两表笔

	c．旋转 S 使其尖端对准欧姆档 $\times 1\mathrm{k}$

	d．旋转 S 使其尖端对准欧姆档 $\times 100$

	e．旋转 S 使其尖端对准交流 $500\UV$ 档，并拔出两表笔

	\tkanswer[2.6cm]{c、a、b、e}。

	根据右图所示指针位置，此被测电阻的阻值约为\tkanswer[1.4cm]{30 k}$\UO$。
	\item
	（多选题）下述关于用多用表欧姆档测电阻的说法中正确的是\xzanswer{AC}

	\fourchoices[answer=AC]
	{测量电阻时如果指针偏转过大，应将选择开关 S 拨至倍率较小的档位，重新调零后测量}
	{测量电阻时，如果红、黑表笔分别插在负、正插孔，则会影响测量结果}
	{测量电路中的某个电阻，应该把该电阻与电路断开}
	{测量阻值不同的电阻时都必须重新调零}
\end{enumerate}
%% answer:
\jdanswer{
\begin{enumerate}
	\item
	c、a、b、e，30 k
	\item
	AC
\end{enumerate}
}
%% memo:
\memoanswer{
（1）题目中要求测量几十 $\UkO$ 的电阻，所以应该选择 $\times 1\mathrm{k}$ 的欧姆挡，再进行欧姆调零，然后进行测量，最后将 S 旋至指向交流 $500\UV$ 挡，拔出表笔，所以顺序为 c、a、b、e。由欧姆表的读数原理，可以得出欧姆表的阻值为 $30\UkO$。

（2）指针偏角过大，说明电阻较小，所以应选择倍率较小的挡位，故 A 正确；电阻的阻值与电流的流向无关，所以红黑表笔接反，不影响测量结果，故 B 错误；不将被测电阻与外电路断开，流经欧姆表的电流会受影响，影响读数甚至烧毁器材，故 C 正确；只要不换挡位，就不用重新调零，故 D 错误。
}

\item
%% number: 16
%% paperName: 2009年上海高考第 16 题 6 分
%% typeId: 4
%% type: 实验题
%% chapter: 电路
%% point: 电动势
%% method: 无
%% score: 6
%% degree: 650
%% duplicateId: 0
%% body:
如图为伏打电池示意图，由于化学反应，在 A、B 两电极附近产生了很薄的两个带电接触层 a、b。
\begin{enumerate}
	\item
	（多选题）沿电流方向绕电路一周，非静电力做功的区域是\xzanswer{BD}
	\fourchoices[answer=BD]
	{$R$}
	{b}
	{$r$}
	{a}
	\item
	在如图所示回路的各区域内，电势升高的总和等于电源的\tkanswer[1.6cm]{电动势}。
\end{enumerate}
\onepicture[w=3.6cm, align=right]{\includesvg{figs/16}}
%% answer:
\jdanswer{
\begin{enumerate}
	\item
	BD
	\item
	电动势
\end{enumerate}
}
%% memo:
\memoanswer{
非静电力做功的区域在电池的内部，但不包括内阻，故 B、D 正确；由电源的工作原理得，电势升高的总和等于电源的电动势。
}

\item
%% number: 17
%% paperName: 2009年上海高考第 17 题 6 分
%% typeId: 4
%% type: 实验题
%% chapter: 运动和力
%% point: 验证牛顿第二定律
%% method: 无
%% score: 6
%% degree: 450
%% duplicateId: 0
%% body:
如图为“用 DIS（位移传感器、数据采集器、计算机）研究加速度和力的关系”的实验装置。
\begin{enumerate}
	\item
	在该实验中必须采用控制变量法，应保持\tkanswer[2.4cm]{小车的总质量}不变，用钩码所受的重力作为\tkanswer[2.4cm]{小车所受外力}，用 DIS 测小车的加速度。
	\item
	改变所挂钩码的数量，多次重复测量。在某次实验中根据测得的多组数据可画出 $a-F$ 关系图线（如图所示）。

	①分析此图线的 OA 段可得出的实验结论是\tkanswer[4.4cm]{在质量不变的条件下，加速度与外力成正比}。

	②（单选题）此图线的 AB 段明显偏离直线，造成此误差的主要原因是\xzanswer{C}

	\fourchoices[answer=C]
	{小车与轨道之间存在摩擦}
	{导轨保持了水平状态}
	{所挂钩码的总质量太大}
	{所用小车的质量太大}
\end{enumerate}
\twopicture[labelA=2009上海17a, labelB=2009上海17b, hA=2.6cm, hB=2.4cm, align=right]
{\includesvg{figs/17a}}
{\includesvg{figs/17b}}
%% answer:
\jdanswer{
\begin{enumerate}
	\item
	小车的总质量，小车所受外力
	\item
	①在质量不变的条件下，加速度与外力成正比；②C
\end{enumerate}
}
%% memo:
\memoanswer{
本题要求研究加速度与力的关系，所以应保持小车的总质量不变，钩码的重力作为合外力。由图象可以得出实验结论：在质量不变的条件下，加速度与外力成正比。本实验的条件是钩码的总质量必须远小于小车的总质量，随着钩码质量的不断增加导致了不满足实验条件，造成误差。
}

\item
%% number: 18
%% paperName: 2009年上海高考第 18 题 6 分
%% typeId: 4
%% type: 实验题
%% chapter: 曲线运动
%% point: 平抛运动
%% method: 无
%% score: 6
%% degree: 450
%% duplicateId: 0
%% body:
利用图~\ref{2009上海18a} 实验可粗略测量人吹气产生的压强。两端开口的细玻璃管水平放置，管内塞有潮湿小棉球，实验者从玻璃管的一端 A 吹气，棉球从另一端 B 飞出，测得玻璃管内部截面积 $S$，距地面高度 $h$，棉球质量 $m$，开始时的静止位置与管口 B 的距离 $x$，落地点 C 与管口 B 的水平距离 $l$。然后多次改变 $x$，测出对应的 $l$，画出 $l^{2}-x$ 关系图线，如图~\ref{2009上海18b} 所示，并由此得出相应的斜率 $k$。
\begin{enumerate}
	\item
	若不计棉球在空中运动时的空气阻力，根据以上测得的物理量可得，棉球从 B 端飞出的速度 $v_{0} =$\tkanswer[1.8cm]{$l\sqrt{\frac{g}{2h}}$}。
	\item
	假设实验者吹气能保持玻璃管内气体压强始终为恒定值，不计棉球与管壁的摩擦，重力加速度 $g$，大气压强 $p_{0}$ 均为已知，利用图~\ref{2009上海18b} 中拟合直线的斜率 $k$ 可得，管内气体压强 $p =$\tkanswer[2.4cm]{$p_{0} + \frac{kmg}{4Sh}$}。
	\item
	考虑到实验时棉球与管壁间有摩擦，则（2）中得到的 $p$ 与实际压强相比\tkanswer[1.0cm]{偏小}（填：偏大、偏小）。
\end{enumerate}
\twopicture[labelA=2009上海18a, labelB=2009上海18b, hA=2.6cm, hB=2.6cm, align=right]
{figs/18a.png}
{figs/18b.png}
%% answer:
\jdanswer{
\begin{enumerate}
	\item
	$l\sqrt{\frac{g}{2h}}$
	\item
	$p_{0} + \frac{kmg}{4Sh}$
	\item
	偏小
\end{enumerate}
}
%% memo:
\memoanswer{
由平抛运动的规律得 $l = v_{0}t$，$h = \frac{1}{2}gt^{2}$，可得 $v_{0} = l\sqrt{\frac{g}{2h}}$。由动能定理得 $(p - p_{0})Sx = \frac{1}{2}mv_{0}^{2}$，并结合 $\frac{l^{2}}{x} = k$，解得 $p = p_{0} + \frac{kmg}{4Sh}$。如果考虑摩擦力，摩擦力要做负功，由 $(p - p_{0})Sx - F_{f}x = \frac{1}{2}mv_{0}^{2}$，可以看出 $p$ 与实际压强相比偏小。
}

\item
%% number: 19
%% paperName: 2009年上海高考第 19 题 4 分
%% typeId: 4
%% type: 实验题
%% chapter: 光学
%% point: 光的电磁说
%% method: 无
%% score: 4
%% degree: 450
%% duplicateId: 0
%% body:
光强传感器对接收到的光信号会产生衰减，且对于不同波长的光衰减程度不同，可以用 $\varphi$ 表示衰减程度，其定义为输出强度与输入强度之比，$\varphi = I_{\text{出}}/I_{\text{入}}$，右图表示 $\varphi$ 与波长 $\lambda$ 之间的关系。当用此传感器分别接收 A、B 两束光时，传感器的输出强度正好相同，已知 A 光的波长 $\lambda_{A} = 625\Unm$，B 光由 $\lambda_{B1} = 605\Unm$ 和 $\lambda_{B2} = 665\Unm$ 两种单色光组成，且这两种单色光的强度之比 $I_{B1}:I_{B2} = 2:3$，由图可知 $\varphi_{A} =$\tkanswer[1.0cm]{0.35}；A 光强度与 B 光强度之比 $I_{A}:I_{B} =$\tkanswer[1.2cm]{27:35}。
\onepicture[h=4.2cm, align=right]{\includesvg{figs/19}}
%% answer: 0.35，27:35
%% memo:
\memoanswer{
（1）由图像看出：当横坐标为 6.25 时，$\varphi_{A} = 0.35$。

（2）当横坐标为 6.05 时，$\varphi_{1} = 0.6$，横坐标为 6.65 时，$\varphi_{2} = 0.05$，由于两种单色光的强度之比 $I_{B1}:I_{B2} = 2:3$，因此 $\varphi_{B} = \frac{2}{5}\times 0.6 + \frac{3}{5}\times 0.05 = 0.27$。由于传感器的输出强度相同，因此 $I_{A}:I_{B} = \varphi_{A}:\varphi_{B} = 27:35$。
}

\item
%% number: 20
%% paperName: 2009年上海高考第 20 题 10 分
%% typeId: 6
%% type: 计算题
%% chapter: 功和能
%% point: 交通工具启动
%% method: 无
%% score: 10
%% degree: 650
%% duplicateId: 0
%% body:
质量为 $5\times10^{3}\Ukg$ 的汽车在 $t = 0$ 时刻速度 $v_{0} = 10\Ums$，随后以 $P = 6\times10^{4}\UW$ 的额定功率沿平直公路继续前进，经 $72\Us$ 达到最大速度，设汽车受恒定阻力，其大小为 $2.5\times10^{3}\UN$。求：
\begin{enumerate}
	\item
	汽车的最大速度 $v_{m}$；
	\item
	汽车在 $72\Us$ 内经过的路程 $s$。
\end{enumerate}
%% answer:
\jdanswer{
\begin{enumerate}
	\item
	$v_{m} = 24\Ums$
	\item
	$s = 1252\Um$
\end{enumerate}
}
%% memo:
\memoanswer{
（1）汽车达到最大速度时，牵引力等于阻力，此时
$P = fv_{m}$
解得 $v_{m} = \frac{P}{f} = \frac{6 \times 10^{4}}{2.5 \times 10^{3}}\Ums = 24\Ums$。

（2）由动能定理得
$Pt - fs = \frac{1}{2}mv_{m}^{2} - \frac{1}{2}mv_{0}^{2}$
$s = \frac{2Pt - mv_{m}^{2} + mv_{0}^{2}}{2f} = 1252\Um$。
}

\item
%% number: 21
%% paperName: 2009年上海高考第 21 题 12 分
%% typeId: 6
%% type: 计算题
%% chapter: 气体性质
%% point: 玻意耳定律
%% method: 无
%% score: 12
%% degree: 450
%% duplicateId: 0
%% body:
如图，粗细均匀的弯曲玻璃管 A、B 两端开口，管内有一段水银柱，右管内气体柱长为 $39\Ucm$，中管内水银面与管口 A 之间气体柱长为 $40\Ucm$。先将 B 封闭，再将左管竖直插入水银槽中，设整个过程温度不变，稳定后右管内水银面比中管内水银面高 $2\Ucm$，求：
\begin{enumerate}
	\item
	稳定后右管内的气体压强 $p$；
	\item
	左管 A 端插入水银槽的深度 $h$。（大气压强 $p_{0} = 76\UcmHg$）
\end{enumerate}
\onepicture[h=3.4cm, align=right]{\includesvg{figs/21}}
%% answer:
\jdanswer{
\begin{enumerate}
	\item
	$p = 78\UcmHg$
	\item
	$h = 7\Ucm$
\end{enumerate}
}
%% memo:
\memoanswer{
（1）插入水银槽后右管内气体，由等温变化得
$p_{0}l_{0} = p\left(l_{0} - \frac{\Delta h}{2}\right)$，
解得 $p = 76\times\frac{39}{38}\UcmHg = 78\UcmHg$。

\onepicture[h=3.6cm, align=right]{\includesvg{figs/21_an}}

（2）插入水银槽后左管压强为
$p' = p + \rho g\Delta h = 80\UcmHg$，
左管内外水银面高度差为
$h_{1} = \frac{p' - p_{0}}{\rho g} = 4\Ucm$，
中、左管内气体，由等温变化得
$p_{0}l = p'l'$
解得 $l' = 38\Ucm$，
左管插入水银槽深度为
$h = l + \frac{\Delta h}{2} - l' + h_{1} = 7\Ucm$。
}

\item
%% number: 22
%% paperName: 2009年上海高考第 22 题 12 分
%% typeId: 6
%% type: 计算题
%% chapter: 运动和力
%% point: 变加速直线运动
%% method: 无
%% score: 12
%% degree: 650
%% duplicateId: 0
%% body:
如图~\ref{2009上海22a}，质量 $m = 1\Ukg$ 的物体沿倾角 $\theta = 37^{\circ}$ 的固定粗糙斜面由静止开始向下运动，风对物体的作用力沿水平方向向右，其大小与风速 $v$ 成正比，比例系数用 $k$ 表示，物体加速度 $a$ 与风速 $v$ 的关系如图~\ref{2009上海22b} 所示。求：
\begin{enumerate}
	\item
	物体与斜面间的动摩擦因数 $\mu$；
	\item
	比例系数 $k$。
\end{enumerate}
\twopicture[labelA=2009上海22a, labelB=2009上海22b, hA=3.0cm, hB=2.6cm, align=right]
{figs/22a.png}
{figs/22b.png}
%% answer:
\jdanswer{
\begin{enumerate}
	\item
	$\mu = 0.25$
	\item
	$k = 0.84\Ukgs$
\end{enumerate}
}
%% memo:
\memoanswer{
（1）对初始时刻有 $v = 0$，$a_{0} = 4\Umsq$。
由牛顿第二定律得
$mg\sin\theta - \mu mg\cos\theta = ma_{0}$，
解得 $\mu = \frac{g\sin\theta - a_{0}}{g\cos\theta} = \frac{6 - 4}{8} = 0.25$。

（2）对末时刻有 $v = 5\Ums$，$a = 0$。
由受力分析得
$mg\sin\theta - \mu N - kv\cos\theta = 0$，
$N = mg\cos\theta + kv\sin\theta$，
$mg(\sin\theta - \mu\cos\theta) - kv(\mu\sin\theta + \cos\theta) = 0$，
解得 $k = \frac{mg(\sin\theta - \mu\cos\theta)}{v(\mu\sin\theta + \cos\theta)} = 0.84\Ukgs$。
}

\item
%% number: 23
%% paperName: 2009年上海高考第 23 题 12 分
%% typeId: 6
%% type: 计算题
%% chapter: 电场
%% point: 带电粒子的圆周运动
%% method: 无
%% score: 12
%% degree: 450
%% duplicateId: 0
%% body:
如图，质量均为 $m$ 的两个小球 A、B 固定在弯成 $120^{\circ}$ 角的绝缘轻杆两端，OA 和 OB 的长度均为 $l$，可绕过 O 点且与纸面垂直的水平轴无摩擦转动，空气阻力不计。设 A 球带正电，B 球带负电，电量均为 $q$，处在竖直向下的匀强电场中。开始时，杆 OB 与竖直方向的夹角 $\theta_{0} = 60^{\circ}$，由静止释放，摆动到 $\theta = 90^{\circ}$ 的位置时，系统处于平衡状态，求：
\begin{enumerate}
	\item
	匀强电场的场强大小 $E$；
	\item
	系统由初位置运动到平衡位置，重力做的功 $W_{g}$ 和静电力做的功 $W_{e}$；
	\item
	B 球在摆动到平衡位置时速度的大小 $v$。
\end{enumerate}
\onepicture[w=6.0cm, align=right]{\includesvg{figs/23}}
%% answer:
\jdanswer{
\begin{enumerate}
	\item
	$E = \frac{mg}{3q}$
	\item
	$W_{g} = \left(\frac{\sqrt{3}}{2} - 1\right)mgl$，$W_{e} = \frac{\sqrt{3}}{6}mgl$
	\item
	$v = \sqrt{\left(\frac{2\sqrt{3}}{3} - 1\right)gl}$
\end{enumerate}
}
%% memo:
\memoanswer{
（1）力矩平衡时：
$(mg - qE)l\sin 90^{\circ} = (mg + qE)l\sin(120^{\circ} - 90^{\circ})$，
即 $mg - qE = \frac{1}{2}(mg + qE)$，
解得 $E = \frac{mg}{3q}$。

（2）重力做功为：
$W_{g} = mgl(\cos 30^{\circ} - \cos 60^{\circ}) - mgl\cos 60^{\circ} = \left(\frac{\sqrt{3}}{2} - 1\right)mgl$，
静电力做功为：
$W_{e} = qEl(\cos 30^{\circ} - \cos 60^{\circ}) + qEl\cos 60^{\circ} = \frac{\sqrt{3}}{6}mgl$。

（3）由动能定理知，系统动能的改变量为
$\Delta E_{k} = W_{g} + W_{e} = \left(\frac{2\sqrt{3}}{3} - 1\right)mgl$，
小球 B 的速度为
$v = \sqrt{\frac{\Delta E_{k}}{m}} = \sqrt{\left(\frac{2\sqrt{3}}{3} - 1\right)gl}$。
}

\item
%% number: 24
%% paperName: 2009年上海高考第 24 题 14 分
%% typeId: 6
%% type: 计算题
%% chapter: 电磁感应
%% point: 电磁感应中的匀变速
%% method: 无
%% score: 14
%% degree: 250
%% duplicateId: 0
%% body:
如图，光滑的平行金属导轨水平放置，电阻不计，导轨间距为 $l$，左侧接一阻值为 $R$ 的电阻。区域 cdef 内存在垂直轨道平面向下的有界匀强磁场，磁场宽度为 $s$。一质量为 $m$，电阻为 $r$ 的金属棒 MN 置于导轨上，与导轨垂直且接触良好，受到 $F = 0.5v + 0.4$（单位为 $\UN$，$v$ 为金属棒运动速度）的水平力作用，从磁场的左边界由静止开始运动，测得电阻两端电压随时间均匀增大。（已知 $l = 1\Um$，$m = 1\Ukg$，$R = 0.3\UO$，$r = 0.2\UO$，$s = 1\Um$）
\begin{enumerate}
	\item
	分析并说明该金属棒在磁场中做何种运动；
	\item
	求磁感应强度 $B$ 的大小；
	\item
	若撤去外力后棒的速度 $v$ 随位移 $x$ 的变化规律满足 $v = v_{0} - \frac{B^{2}l^{2}}{m(R + r)}x$，且棒在运动到 ef 处时恰好静止，则外力 $F$ 作用的时间为多少？
	\item
	若在棒未出磁场区域时撤去外力，画出棒在整个运动过程中速度随位移的变化所对应的各种可能的图线。
\end{enumerate}
\onepicture[w=4.6cm, align=right]{\includesvg{figs/24}}
%% answer:
\jdanswer{
\begin{enumerate}
	\item
	金属棒做匀加速运动
	\item
	$B = 0.5\UT$
	\item
	$t = 1\Us$
	\item
	可能图线如解析图所示
\end{enumerate}
}
%% memo:
\memoanswer{
（1）金属棒做匀加速运动，$R$ 两端电压 $U \propto I \propto E \propto v$，$U$ 随时间均匀增大，即 $v$ 随时间均匀增大，因此加速度为恒量。

（2）由牛顿第二定律得
$F - \frac{B^{2}l^{2}v}{R + r} = ma$，
以 $F = 0.5v + 0.4$ 代入得
$\left(0.5 - \frac{B^{2}l^{2}}{R + r}\right)v + 0.4 = a$。
$a$ 与 $v$ 无关，所以 $a = 0.4\Umsq$。则
$\left(0.5 - \frac{B^{2}l^{2}}{R + r}\right) = 0$，
代入数据得 $B = 0.5\UT$。

（3）由运动学公式知
$x_{1} = \frac{1}{2}at^{2}$，
又 $v_{0} = \frac{B^{2}l^{2}}{m(R + r)}x_{2} = at$，
由 $x_{1} + x_{2} = s$，有
$\frac{1}{2}at^{2} + \frac{m(R + r)}{B^{2}l^{2}}at = s$，
代入数据得 $0.2t^{2} + 0.8t - 1 = 0$，
解得 $t = 1\Us$。

（4）可能图线如下：

\onepicture[w=9.0cm]{\includesvg{figs/24_an}}
}

\end{enumerate}

\end{document}
